% from ega4-2-fr.tex
\oldpage[IV-2]{167}
I) Le corps des fractions $K$ de $A$ est de caractéristique 0. --- On sait
alors (0, 19.8.8) qu'il existe un anneau de valuation discrète complet $C$
et un sous-anneau $B$ de $A$ tel que $A$ soit une $B$-algèbre finie et que
$B$ soit isomorphe à un anneau de séries formelles
$C[[T_1,\ldots,T_n]]$. Comme $C$ est régulier (II, 7.1.6), il en est de
même de $B$ (0, 17.3.8); en outre, le corps des fractions $L$ de $B$ étant
de caractéristique 0, $K$ est une extension séparable finie de $L$, donc on
peut appliquer (6.12.4.1) à $B$, et cela prouve alors la proposition.

II) Le corps des fractions $K$ de $A$ est de caractéristique $p>0$. ---
Alors $A$ contient le corps premier $\mathbf F_p$, et le théorème a été
prouvé dans ce cas (0, 22.7.6).
\end{proof}

\begin{corollary}[6.12.8]
\label{IV.2.s6.12.8-fr}
Soit $A$ un anneau local noethérien complet. Alors, pour tout préschéma $X$
localement de type fini sur $A$, $\operatorname{Reg}(X)$ est ouvert dans
$X$.
\end{corollary}

\begin{proof}
Vérifions la condition c) de (6.12.4). Si $\mathfrak p$ est premier dans
$A$, $A/\mathfrak p$ est encore un anneau local noethérien complet; si $K'$
est une extension finie du corps des fractions $K$ de $A/\mathfrak p$, $K'$
est corps des fractions d'une sous-$A$-algèbre finie $A'$ de $K'$, engendrée
par un système de générateurs de $K'$ sur $K$, entiers sur $A$. On sait
alors que $A'$ est un anneau semi-local complet (Bourbaki, Alg. comm.,
chap. IV, § 2, n$^o$ 5, cor. 3 de la prop. 9), donc produit d'anneaux locaux
complets, et comme $A'$ est intègre, c'est un anneau local complet; en vertu
de (6.12.7), si $X'=\operatorname{Spec}(A')$,
$\operatorname{Reg}(X')$ est ouvert et non vide, d'où la conclusion.
\end{proof}

\begin{proposition}[6.12.9]
\label{IV.2.s6.12.9-fr}
Soit $X$ un préschéma localement noethérien tel que $\operatorname{Reg}(X)$
soit ouvert; alors, pour tout $n\geq0$, l'ensemble $U_{R_n}(X)$ des points
$x\in X$ où $X$ vérifie $(R_n)$ est ouvert.
\end{proposition}

\begin{proof}
En effet (5.8.2), $U_{R_n}(X)$ est réunion de $\operatorname{Reg}(X)$ et de
l'ensemble des $x\in\operatorname{Sing}(X)$ tels que les points génériques
$z_i$ des composantes irréductibles de $\operatorname{Sing}(X)$ contenant
$x$ vérifient la relation $\dim(\mathcal O_{X,z_i})\geq n+1$ (puisque l'on a
$z_i\notin\operatorname{Reg}(X)$); autrement dit, ces points
$x\in\operatorname{Sing}(X)$ sont ceux pour lesquels on a
$\operatorname{codim}_x(\operatorname{Sing}(X),X)\geq n+1$ (5.1.2); comme
la fonction $x\mapsto\operatorname{codim}_x(\operatorname{Sing}(X),X)$ est
semi-continue intérieurement (0, 14.2.6), l'ensemble de ces points est
ouvert dans $\operatorname{Sing}(X)$, ce qui achève la démonstration.
\end{proof}

Notons aussi le résultat élémentaire suivant:

\begin{proposition}[6.12.10]
\label{IV.2.s6.12.10-fr}
Soit $X$ un préschéma localement noethérien. L'ensemble $U_{R_0}(X)$ est
ouvert dans $X$. Pour que $U_{R_1}(X)$ soit ouvert dans $X$, il faut et il
suffit que tout point maximal $x\in X$ tel que $x\in U_{R_1}(X)$ (ce qui
signifie que $X$ est réduit en $x$) soit intérieur à $U_{R_1}(X)$.
\end{proposition}

\begin{proof}
Par définition (5.8.2), $U_{R_0}(X)$ est l'ensemble
$X-\bigcup_\alpha X'_\alpha$, où $X'_\alpha$ parcourt l'ensemble des
composantes irréductibles de $X$ telles que $X$ ne soit pas réduit en leur
point générique; comme l'ensemble des composantes irréductibles de $X$ est
localement fini, $U_{R_0}(X)$ est ouvert. En ce qui concerne $U_{R_1}(X)$,
la condition de l'énoncé étant trivialement nécessaire, prouvons qu'elle est
suffisante. Soient $X''_\beta$ les composantes irréductibles de $X$ telles
qu'en leur point générique $x''_\beta$ le préschéma $X$ soit réduit; par
hypothèse il existe un ouvert $U\subset U_{R_1}(X)$ contenant tous les
$x''_\beta$. Désignons alors par $Z_\lambda$ les composantes irréductibles
de l'ensemble fermé $Z=X-U$ dont le point générique $z_\lambda$ est tel que
$\dim(\mathcal O_{X,z_\lambda})\leq1$ et que
$\mathcal O_{X,z_\lambda}$ soit non régulier. Aucun point appartenant à
l'un des $Z_\lambda$ ne peut appartenir
